\documentclass[11pt]{article}
\usepackage[T1]{fontenc}
\usepackage{lmodern}
\usepackage{amsmath,amssymb,amsthm}
\usepackage{graphicx}
\usepackage[margin=1in]{geometry}
\usepackage[colorlinks=true,linkcolor=blue,citecolor=blue,urlcolor=blue]{hyperref}

\newtheorem{theorem}{Theorem}
\newtheorem{corollary}{Corollary}
\newtheorem{remark}{Remark}
\newcommand{\T}{\mathbb T}
\newcommand{\Z}{\mathbb Z}

\title{A Harmonic-Near-Symmetric Order Obstruction\\
at the Critical Sobolev Index}
\author{Run \texttt{fa\_banach\_001}, agent \texttt{agent\_lane\_11}}
\date{17 August 2026}

\begin{document}
\maketitle

\section*{Source Question and Status}

Ivanov and Kalton study when a Riesz basis of exponentials in
$L^2(-\pi,\pi)$ remains a basis in Sobolev spaces.  After proving their
unconditional-basis classification, they ask whether the exponentials can be
a \emph{conditional} basis of $H^s$, under a suitable ordering, throughout
the intermediate range $s_0\leq s\leq s_\Lambda$ \cite[p.~14]{IK}.

\begin{center}
\includegraphics[width=0.96\linewidth]{figures/open_problem_crop.png}
\end{center}

This packet gives a rigorous negative answer for a broad class of orders in
the canonical endpoint model $\Lambda=\Z$, where $s_0=s_\Lambda=1/2$.
The full question for arbitrary, highly nonlocal permutations remains open.

\section*{Definitions}

Write $I=(-\pi,\pi)$.  We use the equivalent Slobodeckij norm
\[
 \|g\|_{H^{1/2}(I)}^2\asymp \|g\|_{L^2(I)}^2+[g]_{1/2,I}^2,
 \qquad
 [g]_{1/2,I}^2=
 \int_I\!\int_I\frac{|g(x)-g(y)|^2}{|x-y|^2}\,dx\,dy.
\]
Let $\pi_0:\mathbb N\to\Z$ be an enumeration and
$A_k=\{\pi_0(1),\ldots,\pi_0(k)\}$.  We call the order
\emph{harmonic-near-symmetric along a subsequence} if there are
$k_j,N_j\to\infty$ such that
\begin{equation}\label{eq:harmonic}
 \sum_{\substack{n\in A_{k_j}\mathbin\triangle[-N_j,N_j]\\n\ne0}}
 \frac1{|n|}\longrightarrow0.
\end{equation}

\section*{Proof Intuition}

The smooth interval function $f(x)=x$ has Fourier coefficients of size
$1/|n|$.  Its symmetric Fourier sums converge in every subcritical
$H^s(I)$, but at $s=1/2$ their periodic jump is resolved in a boundary layer
of width $1/N$.  After rescaling that layer, the error tends to a nonconstant
profile.  The scale invariance of the critical Slobodeckij seminorm turns the
profile into a fixed positive lower bound.  Changing the retained frequencies
by a set of harmonic mass $o(1)$ changes the sum by $o(1)$ in $H^{1/2}(I)$,
so every order satisfying \eqref{eq:harmonic} inherits the defect.

\section*{Main Result}

\begin{theorem}\label{thm:main}
If an enumeration $\pi_0:\mathbb N\to\Z$ is harmonic-near-symmetric along a
subsequence, then the ordered system
\[
  \bigl(e^{i\pi_0(k)x}\bigr)_{k\geq1}
\]
is not a Schauder basis of $H^{1/2}(-\pi,\pi)$.
\end{theorem}

\begin{proof}
Let $f(x)=x$.  This function belongs to $H^{1/2}(I)$, and with the
normalization
\(
 \widehat f(n)=(2\pi)^{-1}\int_{-\pi}^{\pi}f(x)e^{-inx}\,dx
\)
one has
\[
 \widehat f(0)=0,
 \qquad
 \widehat f(n)=\frac{i(-1)^n}{n}\quad(n\ne0).
\]
Its symmetric Fourier sum is therefore
\[
 S_Nf(x)=2\sum_{n=1}^N\frac{(-1)^{n+1}\sin(nx)}n.
\]

We first prove that these sums do not converge to $f$ in $H^{1/2}(I)$.
At the boundary scale $x=\pi-u/N$ their errors satisfy
\begin{align*}
 r_N(\pi-u/N)
 &:=f(\pi-u/N)-S_Nf(\pi-u/N)\\
 &=\pi-\frac uN-2\sum_{n=1}^N\frac{\sin(nu/N)}n.
\end{align*}
For $u$ in any fixed compact subset of $[0,\infty)$,
\[
 \sum_{n=1}^N\frac{\sin(nu/N)}n
 =\frac1N\sum_{n=1}^N h(u,n/N),
 \qquad
 h(u,t)=
 \begin{cases}\sin(ut)/t,&t>0,\\u,&t=0,
 \end{cases}
\]
is a Riemann sum converging uniformly to $\int_0^1\sin(ut)/t\,dt$.
Consequently
\begin{equation}\label{eq:profile}
 r_N(\pi-u/N)\longrightarrow
 R(u):=\pi-2\int_0^1\frac{\sin(ut)}t\,dt
\end{equation}
uniformly on compact $u$-intervals.

The profile is nonconstant: $R'(u)=-2\sin(u)/u$ for $u>0$.  Choose two
disjoint compact intervals $U,V\subset(0,\pi)$ on which the values of $R$
are separated by some $\delta>0$.  Uniform convergence in
\eqref{eq:profile} then gives, for all sufficiently large $N$,
\[
 |r_N(\pi-u/N)-r_N(\pi-v/N)|\geq\delta/2
 \qquad(u\in U,\ v\in V).
\]
Restricting the Slobodeckij integral to the corresponding two boundary
strips and changing variables yields
\begin{align*}
 [r_N]_{1/2,I}^2
 &\geq
 \int_U\!\int_V
 \frac{|r_N(\pi-u/N)-r_N(\pi-v/N)|^2}{|u-v|^2}\,du\,dv\\
 &\geq c_{U,V,\delta}>0.
\end{align*}
Thus $\liminf_N\|f-S_Nf\|_{H^{1/2}(I)}>0$.

For a finite set $A\subset\Z$, let
$T_Af=\sum_{n\in A}\widehat f(n)e^{inx}$.  Restriction from the periodic
space $H^{1/2}(\T)$ to $H^{1/2}(I)$ is continuous.  Hence, with
$D=A\mathbin\triangle[-N,N]$,
\begin{align}
 \|(T_A-S_N)f\|_{H^{1/2}(I)}^2
 &\leq C\|(T_A-S_N)f\|_{H^{1/2}(\T)}^2 \notag\\
 &\leq C'\sum_{\substack{n\in D\\n\ne0}}
 (1+n^2)^{1/2}|\widehat f(n)|^2
 \leq C''\sum_{\substack{n\in D\\n\ne0}}\frac1{|n|}.
 \label{eq:perturb}
\end{align}
By \eqref{eq:harmonic},
$\|(T_{A_{k_j}}-S_{N_j})f\|_{H^{1/2}(I)}\to0$.  The boundary lower bound and
the triangle inequality imply
\[
 \liminf_j\|f-T_{A_{k_j}}f\|_{H^{1/2}(I)}>0.
\]

Finally, if the ordered exponentials were a Schauder basis, its expansion
coefficients for $f$ would have to be $\widehat f(n)$.  Indeed, convergence
in $H^{1/2}(I)$ implies convergence in $L^2(I)$, where the integer
exponentials are an orthogonal basis and coefficients are unique.  Its
$k$th basis partial sum would therefore be $T_{A_k}f$, contradicting the
display above.
\end{proof}

\section*{Concrete Orders Ruled Out}

Let $\eta(1)=0$, $\eta(2r)=r$, and $\eta(2r+1)=-r$ be the canonical
symmetric enumeration, and let $p(m)$ be the position occupied by $\eta(m)$
in another enumeration.

\begin{corollary}
The conclusion of Theorem~\ref{thm:main} holds if either
\begin{enumerate}
\item $\sup_m|p(m)-m|<\infty$, or
\item $|p(m)-m|=o(m)$ as $m\to\infty$.
\end{enumerate}
In particular, the symmetric order, every bounded positional perturbation,
and every sublinear positional perturbation fail at $s=1/2$.
\end{corollary}

\begin{proof}
Let $B_M=\{\eta(1),\ldots,\eta(M)\}$ and let $A_M$ be the first $M$
frequencies in the new order.  If $|p(m)-m|\leq C$, membership in $A_M$ and
$B_M$ can differ only when $|m-M|\leq C$; hence their harmonic symmetric
difference has mass $O(C/M)$.  If $|p(m)-m|=o(m)$, fix $\varepsilon>0$.
For all sufficiently large $m$,
$(1-\varepsilon)m\leq p(m)\leq(1+\varepsilon)m$.  Membership can then differ
only for
\[
 \frac{M}{1+\varepsilon}<m<\frac{M}{1-\varepsilon},
\]
apart from finitely many indices, which agree once $M$ is large.  This window
has $O(\varepsilon M)$ indices, and $|\eta(m)|\asymp M$ there, so its harmonic
mass is $O(\varepsilon)$.  Taking $M=2N+1$ gives $B_M=[-N,N]$; first let
$N\to\infty$ and then $\varepsilon\downarrow0$ to obtain
\eqref{eq:harmonic}.
\end{proof}

\section*{Scope, Upgrade Attempts, and Novelty Check}

The theorem is not a full answer to the Ivanov--Kalton question.  An arbitrary
permutation can retain a reservoir of remote frequencies with positive
harmonic mass while delaying frequencies at every scale, so
\eqref{eq:perturb} need not compare any of its prefixes with a symmetric sum.
The argument also does not cover general noninteger $\Lambda$.

Eight focused upgrade attempts were made before promotion: a false
Bessel-to-Riesz implication was eliminated; uniform minimality was checked;
the lattice endpoint was reduced to exact spectral finite sections; the
boundary profile was proved; it was extended to harmonic-near-symmetric
orders; arbitrary permutations were tested and isolated as the obstruction;
weighted conditional exponential bases were examined as possible positive
models; and regular split-shift examples from the source were stress-tested.
The last computation again showed endpoint growth for natural symmetric
orders but supplied no theorem for exotic permutations.

The run's cheap indexes and a bounded exact-phrase/citation search through
17 August 2026 found no later answer to the exact question.  This theorem
should therefore be reviewed as a new ordering-sensitive partial result, not
as a claim that the source problem is resolved.

\section*{Verification Notes}

The proof is analytic.  The companion script evaluates the limiting profile,
checks finite-$N$ convergence on a fixed grid, and verifies explicit
separation between two profile intervals.  These computations are consistency
checks only and are not used in the proof.

\section*{Human Review Recommendation}

Verify the critical boundary rescaling, the periodic-to-interval norm
comparison in \eqref{eq:perturb}, and the positional-displacement corollary.
Most importantly, preserve the distinction between the proved
harmonic-near-symmetric obstruction and the unresolved arbitrary-order class.

\enlargethispage{2\baselineskip}
\begin{thebibliography}{9}
\small

\bibitem{IK}
S. Ivanov and N. Kalton,
\emph{Interpolation of subspaces and applications to exponential bases in
Sobolev spaces},
arXiv:math/0104130; St. Petersburg Math. J. \textbf{13} (2002), no.~2,
221--239.

\end{thebibliography}

\end{document}
